
\section{Introduction}














Many reinforcement-learning algorithms follow the policy-iteration paradigm, which alternates between 
estimating the action values of a policy 
and improving the policy using those estimates. 
This framework is fundamental because it can learn an optimal policy for an unknown Markov decision process (MDP) using only sampled rewards and transitions.
Classical policy iteration improves the policy once the action-value estimates have converged.
In contrast, optimistic variants improve the policy after partial evaluation.
This reduces computation as it requires fewer evaluation steps between each improvement step.
However, it complicates the analysis because 

small perturbations 

to the value estimates
can drastically affect the policy and alter the subsequent dynamics.


















Convergence guarantees for optimistic policy-iteration algorithms remain limited, even in the tabular setting.
Algorithms that use one-step bootstrapped estimates for evaluation, such as $Q$-learning, converge to optimality under standard asynchronous stochastic-approximation conditions.
When using multi-step bootstrapped estimates, convergence to optimality requires additional structure, such as combining several multi-step approximations~\citep{Harutyunyan2016QCorrections, Munos2016SafeLearning} or using a look-ahead mechanism~\citep{Winnicki2023OnEvaluation}.
Monte Carlo Optimistic Policy Iteration (MC-O-PI), which uses full-episode returns for evaluation, remains poorly understood. Despite its simplicity, the asymptotic behaviour of this foundational algorithm has been an open question for three decades~\citep{Sutton1999OpenLearning}.
Convergence to optimality is guaranteed when the MDP exhibits a feedforward structure \citep{Wang2022OnLearning, Lubars2021OptimisticStructure}
or when the discount factor is smaller than $1/2$ \citep{Delattre2025MarkovAlgorithm}.
The classical guarantees for general finite MDPs without a small-discount restriction require that all state-action pairs are updated at the same frequency \citep{Tsitsiklis2002OnIteration, Chen2018OnProblem, Liu2021OnStarts}.














Updating all pairs as frequently requires that episodes are initialised uniformly over the entire state-action space.
This condition is unrealistic in practice, especially when the state space is large.
However, once the initial state of an episode is selected, uniformly sampling the initial action is usually straightforward.
This paper proves that this weaker condition is sufficient.
Specifically, initial-visit MC-O-PI converges almost surely to the optimal action values when updates are uniform over the actions within each state.
Different states may thus be updated at arbitrary frequencies, provided each receives infinite accumulated learning. The precise sampling, policy-selection and learning-rate conditions are collected in \autoref{asp: standing}.
This result relaxes the state-sampling requirement of the previous convergence guarantees 
and applies in settings where uniform state sampling is not feasible but uniform action sampling is easy.






















The proof of this new theorem, stated as \autoref{thm: convergence mcopi uniform} below, requires new tools. 
Previous work demonstrates convergence by studying the evolution of the value estimates. The arguments rely centrally on a commutativity property of the Bellman operator that fails when updates are not uniform over all state-action pairs.
Instead, we track the evolution of the policy and find that the mean-field dynamics of MC-O-PI generate monotonically improving policies when updates are uniform over the actions within each state.
We then show that stochastic perturbations cannot systematically defeat this improvement: at sufficiently late transitions, there is a fixed positive chance of improving the policy or of keeping the same action-value target forever. Boundedness allows this probability estimate to be made before exit from a fixed bounded set. Repeated opportunities then imply that only finitely many target changes occur, so the value estimates converge to the optimal action values.

This develops the combined stability--ODE method of \citet[Section~5.4]{Kushner2003StochasticApplications} for entries into policy regions that may approach their boundaries. The proof explains how the learning-rate condition allows noise control even when these entry margins vanish.

This paper is organised as follows. Section~2 introduces finite Markov decision
processes and formulates MC-O-PI as a stochastic difference inclusion. Section~3 proves the
main convergence theorem, first for the mean-field dynamics and then for the stochastic
algorithm. Section~4 discusses the practical implications of action-wise uniform updates
and possible extensions to related optimistic methods. 

Unless stated otherwise, operations and relations between functions are elementwise. For instance, given functions $f, g : \mathcal X \to \mathcal Y$ we write
\begin{align*}
&(f + g)(x) = f(x) + g(x), \qquad \forall \, x \in \xset ,
    \\
&f \le g \iff f(x) \le g(x), \qquad \forall \, x \in \xset .
\end{align*}
Moreover, we compute norms as $\|f\| := \sup_{x \in \xset} |f(x)|$ 
and adopt the convention $\inf \varnothing = \infty$.

























































































